\subsection*{3.3 Mainline synthesis: residual scalar, restoring law, and reversal probe}
\addcontentsline{toc}{subsection}{3.3 Mainline synthesis: residual scalar, restoring law, and reversal probe}
The 3.3 compression is now strong enough to be stated directly in the main text. The visible residual sector may be organized around the single transported scalar
\[
  \eta:=\frac13-\beta^2.
\]
In the symmetric readout one has $\eta_{\mathrm{sym}}\approx 3.1603\times10^{-2}$ and therefore
\[
  \eta_{\mathrm{sym}}^2\approx 9.987\times10^{-4},
\]
so a number near $10^{-3}$ is read most naturally as a quadratic residual of the same transported defect sector rather than as an additional primitive constant.

\begin{theorem}[3.3 clean mainline synthesis]
\label{thm:mainline-332}
On an admissible local branch of the transported closure-stable kernel, assume that the visible compensated defect is isotropic and that no new primitive local scalar channel opens beyond the continuum--discrete deformation already seen in $\beta$. Then the late 3.3 structure reduces locally to four linked statements:
\begin{enumerate}[label=(\roman*)]
  \item the visible residual sector is governed by the single scalar $\eta=\frac13-\beta^2$;
  \item the leading Julia-boundary defect is controlled by the same scalar,
  \[
    \mathfrak J = c\,\eta + O(\eta^2);
  \]
  \item the local closure-restoring return law takes the normal form
  \[
    \eta''+\omega^2\eta = O(\eta^2), \qquad \omega^2>0;
  \]
  \item the hybrid reversal probe
  \[
    \mathfrak S_{\mathrm{rev}}=(\mathbb P_{\mathrm B},\mathcal L_{\mathrm L})
  \]
  has only the two local outcomes
  \[
    \Delta\theta\equiv 0 \pmod{2\pi}
    \qquad\text{or}\qquad
    \Delta\theta\equiv \pi \pmod{2\pi},
  \]
  i.e. direct convergence or closure only on the spinorial/M\"obius double cover.
\end{enumerate}
The supporting compressed proofs and the detailed transport lemmas remain collected in Appendix~C.
\end{theorem}

\begin{proof}
Appendix~C proves that the visible defect, the Julia-boundary readout, the local Umkehrwalze, and the reversal-probe dichotomy are not separate mechanisms but four faces of the same transported residual sector. For the main text it is therefore enough to record the resulting normal form here, while the denser proof infrastructure and transport tools remain collected in Appendix~C.
\end{proof}

\subsection*{3.4 First extension: closure potential and global branch law}
\addcontentsline{toc}{subsection}{3.4 First extension: closure potential and global branch law}
The first 3.4 step is to globalize the local Umkehrwalze. Instead of adding more local residual lemmas, we ask whether the autonomous reduced branch equation from Appendix~C is exact on a connected admissible branch interval. If so, the visible residual dynamics is generated by a one-variable closure potential and the reversal probe admits an energy-form reading.

\begin{theorem}[3.4 first globalization under closure exactness]
\label{thm:mainline-340}
Assume the hypotheses of Theorem~\ref{thm:mainline-332}. Assume in addition that on a connected admissible branch interval $I\ni 0$ the reduced equation
\[
  \eta'' = B(\eta)
\]
is exact, i.e. there exists $\mathcal U\in C^2(I)$ with
\[
  B(\eta) = -\mathcal U'(\eta).
\]
Then:
\begin{enumerate}[label=(\roman*)]
  \item the local Umkehrwalze globalizes on $I$ to the branch law
  \[
    \eta''+\partial_\eta\mathcal U(\eta)=0;
  \]
  \item the branch energy
  \[
    \mathcal E_\eta := \frac12(\eta')^2 + \mathcal U(\eta)
  \]
  is conserved along admissible trajectories;
  \item if the branch is closure-restoring, then $\eta=0$ is a strict local minimum of $\mathcal U$ and
  \[
    \mathcal U(\eta) = \frac12 \omega^2 \eta^2 + O(\eta^3),
    \qquad \omega^2>0;
  \]
  \item if, on each admissible regular energy component, no monodromy higher than the spinorial $\mathbb Z_2$-lift opens, then the reversal probe is globally reduced to the same two-way outcome as locally: direct convergence on the visible branch or closure only on the M\"obius double cover.
\end{enumerate}
The supporting exactness and energy arguments are recorded in Appendix~C as the first 3.4 working extension.
\end{theorem}

\begin{proof}
Appendix~C proves the exact-potential reduction under branch-local closure exactness. Differentiating the energy identity gives $d\mathcal E_\eta/ds = \eta'(\eta''+\mathcal U'(\eta))=0$. The local 3.3 restoring law already shows that the exact core is nondegenerate and closure-restoring, so the potential has a strict minimum there with Hessian $\omega^2>0$. The global probe claim is the energy-form lift of the local 3.3 dichotomy: once the hard admissibility filter and the fine residual functional place two probes on the same regular admissible energy component, only direct closure or the unique nontrivial $\mathbb Z_2$ lift remains.
\end{proof}


\subsection*{3.4 Second extension: reversible single-well potential class}
\addcontentsline{toc}{subsection}{3.4 Second extension: reversible single-well potential class}
The next 3.4 step is to constrain the potential itself. Once the visible residual branch admits an exact one-variable closure potential, the admissible spinorial reversal already suggests that the potential cannot be arbitrary: under branch reversibility and uniqueness of the exact core, the only natural class is a symmetric single-well potential whose scalar ambiguity is exhausted by the transported $\mathbb Z_2$-lift.

\begin{theorem}[3.4 reversible single-well restriction]
\label{thm:mainline-341}
Assume the hypotheses of Theorem~\ref{thm:mainline-340}. Assume moreover that on the connected admissible branch interval $I$:
\begin{enumerate}[label=(\alph*)]
  \item the admissible residual branch is invariant under spinorial reversal $R(\eta)=-\eta$;
  \item the reduced branch law commutes with $R$;
  \item the exact compensated core is the unique admissible critical core on $I$, and $\mathcal U$ has no further admissible critical point on $I\setminus\{0\}$.
\end{enumerate}
Normalize $\mathcal U(0)=0$. Then:
\begin{enumerate}[label=(\roman*)]
  \item $B$ is odd and $\mathcal U$ is even on $I$;
  \item the local restoring expansion improves to the parity-restricted form
  \[
    \mathcal U(\eta)=\frac12\omega^2\eta^2+a_4\eta^4+O(\eta^6),
    \qquad
    B(\eta)=-\omega^2\eta-4a_4\eta^3+O(\eta^5);
  \]
  \item $\eta=0$ is the unique admissible well bottom, $\mathcal U(\eta)>0$ for every $\eta\neq 0$ in $I$, and $\mathcal U$ is strictly increasing as a function of $|\eta|$ on each half-branch;
  \item every admissible regular energy shell has at most one positive and one negative turning value, so the scalar global ambiguity is exhausted by the transported pair $\eta\leftrightarrow-\eta$; hence, together with the no-higher-monodromy hypothesis of Theorem~\ref{thm:mainline-340}, the global reversal-probe ambiguity remains exactly the two-way alternative: direct closure or spinorial/M\"obius lift.
\end{enumerate}
The supporting symmetry and monotonicity argument is recorded in Appendix~C as the second 3.4 extension.
\end{theorem}

\begin{proof}
Appendix~C shows that reversibility forces odd/even parity of the reduced branch law and the closure potential once the exact branch law is available. The nondegenerate closure-restoring core then gives the positive quadratic term, while parity removes the odd Taylor coefficients. Uniqueness of the admissible core and absence of further admissible critical points prevent any secondary well or plateau, so the potential is single-well and monotone in $|\eta|$ on each half-branch. Therefore a regular energy shell can contribute at most the symmetric turning pair $\pm\eta_E$, and the remaining global ambiguity is precisely the transported $\mathbb Z_2$ branch pair already singled out by the reversal probe.
\end{proof}

\subsection*{3.4 Third extension: canonical quartic branch family}
\addcontentsline{toc}{subsection}{3.4 Third extension: canonical quartic branch family}
The next 3.4 step is to make the single-well class explicit without overfitting it. Once parity, reversibility, and core uniqueness have already removed all odd terms and all additional admissible minima, the first nontrivial nonlinear correction is forced into the quartic sector. If no further admissible local scalar or independent local scale opens at this order, then the visible branch law collapses to the reversible quartic Duffing family.

\begin{theorem}[3.4 canonical quartic model family]
\label{thm:mainline-342}
Assume the hypotheses of Theorem~\ref{thm:mainline-341}. Assume in addition that on a symmetric admissible branch neighborhood $J\subset I$ of $0$:
\begin{enumerate}[label=(\alph*)]
  \item $\mathcal U$ is $C^6$;
  \item no additional admissible local scalar or independent local scale enters the first nonlinear correction of the exact one-field branch law;
  \item the first nontrivial nonlinear correction is already visible on $J$.
\end{enumerate}
Then there exists $\lambda\ge 0$ and a remainder $R_6(\eta)=O(\eta^6)$ such that on $J$
\[
  \mathcal U(\eta)=\frac12\omega^2\eta^2+\frac{\lambda}{4}\eta^4+R_6(\eta),
\]
and hence
\[
  B(\eta)=-\omega^2\eta-\lambda\eta^3+O(\eta^5).
\]
In the minimal truncation $R_6\equiv 0$, the exact visible branch law is the reversible quartic family
\[
  \eta''+\omega^2\eta+\lambda\eta^3=0,
  \qquad \lambda\ge 0,
\]
with conserved branch energy
\[
  \mathcal E_\eta = \frac12(\eta')^2+\frac12\omega^2\eta^2+\frac{\lambda}{4}\eta^4.
\]
Moreover:
\begin{enumerate}[label=(\roman*)]
  \item $\lambda=0$ is the harmonic core case, while $\lambda>0$ is the genuinely nonlinear single-well branch law;
  \item for every regular energy $E>0$, the turning values are the symmetric pair $\pm A(E)$ with $A(E)>0$ determined by
  \[
    E=\frac12\omega^2A(E)^2+\frac{\lambda}{4}A(E)^4;
  \]
  \item if $\lambda>0$, then $A(E)$ is strictly increasing in $E$, so the global probe can equivalently be read through the energy shell or through the unique positive turning value.
\end{enumerate}
The supporting reduction is recorded in Appendix~C as the third 3.4 extension.
\end{theorem}

\begin{proof}
Theorem~\ref{thm:mainline-341} already forces the potential to be even and single-well, with local expansion
\[
  \mathcal U(\eta)=\frac12\omega^2\eta^2+a_4\eta^4+O(\eta^6).
\]
By the present hypotheses, no independent admissible local scalar or local scale appears at the first nonlinear order, so the quartic coefficient is the only visible nonlinear datum. Write $\lambda:=4a_4$. Since the potential is single-well with strict minimum at $0$, the first admissible nonlinear correction cannot open a secondary critical point on $J$; shrinking $J$ if necessary, this forces $\lambda\ge 0$. The branch law follows by differentiation. In the minimal truncation the energy identity is immediate. The turning values solve $\eta'=0$, hence $E=\mathcal U(\eta)$; by even single-well monotonicity this yields exactly one positive turning value $A(E)$ and its negative partner. Differentiating the turning relation gives monotonicity of $A(E)$, and therefore the regular probe shell can be parameterized either by $E$ or by $A(E)$.
\end{proof}

\subsection*{3.4 Fourth extension: quartic closure coefficient and visible nonlinearity}
\addcontentsline{toc}{subsection}{3.4 Fourth extension: quartic closure coefficient and visible nonlinearity}
The quartic family is still only a kinematic normal form until its nonlinear coefficient is tied back to the closure side. The cleanest compatible statement is that, once the exact branch potential is the visible readout of an even closure geometry on the same admissible branch, the quartic coefficient is not free: it is exactly the transported quartic closure coefficient.

\begin{theorem}[3.4 quartic closure readout]
\label{thm:mainline-343}
Assume the hypotheses of Theorem~\ref{thm:mainline-342}. Assume moreover that on a symmetric admissible neighborhood $J$ of $0$ there exist an even $C^6$ closure scalar $\Gamma_{\mathrm{cl}}$ and a transport factor $\mu>0$ such that
\[
  \mathcal U(\eta)=\mu\bigl(\Gamma_{\mathrm{cl}}(\eta)-\Gamma_{\mathrm{cl}}(0)\bigr)+O(\eta^6)
\]
and
\[
  \Gamma_{\mathrm{cl}}(\eta)
  =\Gamma_{\mathrm{cl}}(0)+\frac{c_2}{2}\eta^2+\frac{c_4}{4}\eta^4+O(\eta^6),
  \qquad c_2>0,\quad c_4\ge 0.
\]
Then the visible quadratic and quartic branch coefficients satisfy
\[
  \omega^2=\mu c_2,
  \qquad
  \lambda=\mu c_4.
\]
In particular, the visible quartic coefficient is the transported quartic closure coefficient, and
\[
  \frac{\lambda}{\omega^2}=\frac{c_4}{c_2}.
\]
Consequently:
\begin{enumerate}[label=(\roman*)]
  \item $c_4=0$ if and only if $\lambda=0$, i.e. the visible branch remains harmonic at quartic order;
  \item $c_4>0$ if and only if $\lambda>0$, i.e. the visible branch carries a genuine nonlinear single-well correction;
  \item the first nonlinear visible datum is geometrically fixed by closure and is not an additional free branch parameter.
\end{enumerate}
The supporting coefficient-matching argument is recorded in Appendix~C as the fourth 3.4 extension.
\end{theorem}

\begin{proof}
Insert the closure expansion into the transport relation. One obtains
\[
  \mathcal U(\eta)
  =\frac{\mu c_2}{2}\eta^2+\frac{\mu c_4}{4}\eta^4+O(\eta^6).
\]
Comparing with Theorem~\ref{thm:mainline-342},
\[
  \mathcal U(\eta)=\frac12\omega^2\eta^2+\frac{\lambda}{4}\eta^4+O(\eta^6),
\]
forces $\omega^2=\mu c_2$ and $\lambda=\mu c_4$. Since $\mu>0$, the vanishing and positivity statements for $c_4$ and $\lambda$ are equivalent. The ratio identity is immediate. Thus the quartic visible correction is exactly the transported quartic closure coefficient and introduces no new independent nonlinear branch datum.
\end{proof}



\subsection*{3.4 Fifth extension: nonlinear period drift and probe readout}
\addcontentsline{toc}{subsection}{3.4 Fifth extension: nonlinear period drift and probe readout}
Once the quartic coefficient is tied to closure, the next question is no longer merely which family is admissible, but how that nonlinear datum becomes visible to the probe. For the reversible quartic family, the first nonlinear timing defect is already controlled by the same coefficient. In other words: the first deviation from local isochrony is not a new parameter, but the operational readout of the quartic closure term itself.

\begin{theorem}[3.4 nonlinear period-drift readout]
\label{thm:mainline-344}
Assume the hypotheses of Theorem~\ref{thm:mainline-343}. Consider the quartic exact branch law
\[
  \eta''+\omega^2\eta+\lambda\eta^3=0,
  \qquad \omega^2>0,\quad \lambda\ge 0,
\]
on a regular small-energy shell. Let $T(E)$ denote the minimal positive period of the corresponding nontrivial periodic branch solution with conserved energy $E>0$, and let $A(E)>0$ be its positive turning value. Then, as $E\downarrow 0$,
\[
  T(E)=\frac{2\pi}{\omega}\Bigl(1-\frac{3\lambda}{4\omega^4}E+O(E^2)\Bigr),
\]
and equivalently, as $A\downarrow 0$,
\[
  T(A)=\frac{2\pi}{\omega}\Bigl(1-\frac{3\lambda}{8\omega^2}A^2+O(A^4)\Bigr).
\]
Using Theorem~\ref{thm:mainline-343}, this may be rewritten as
\[
  T(E)=\frac{2\pi}{\omega}\Bigl(1-\frac{3c_4}{4\mu c_2^2}E+O(E^2)\Bigr).
\]
Consequently:
\begin{enumerate}[label=(\roman*)]
  \item $\lambda=0$ if and only if $T(E)=\tfrac{2\pi}{\omega}+O(E^2)$, i.e. the branch is locally isochronous to first nonlinear order;
  \item $\lambda>0$ implies $T'(E)<0$ for all sufficiently small $E>0$, so the period decreases monotonically with energy at first nonlinear order;
  \item the first nonlinear timing drift is therefore an operational readout of the same quartic closure coefficient $c_4$, and introduces no further visible scalar.
\end{enumerate}
The supporting perturbative expansion is recorded in Appendix~C as the fifth 3.4 extension.
\end{theorem}

\begin{proof}
Write the periodic solution in the form $\eta(s)=A\cos(\Omega s)+O(A^3)$ and insert it into
\[
  \eta''+\omega^2\eta+\lambda\eta^3=0.
\]
A standard Poincar\'e--Lindstedt balance cancels the resonant $\cos(\Omega s)$ term and yields
\[
  \Omega(A)=\omega+\frac{3\lambda}{8\omega}A^2+O(A^4).
\]
Hence
\[
  T(A)=\frac{2\pi}{\Omega(A)}
  =\frac{2\pi}{\omega}\Bigl(1-\frac{3\lambda}{8\omega^2}A^2+O(A^4)\Bigr).
\]
On the same small shell,
\[
  E=\frac12\omega^2A^2+\frac{\lambda}{4}A^4+O(A^6),
\]
so $A^2=\tfrac{2E}{\omega^2}+O(E^2)$. Substituting gives
\[
  T(E)=\frac{2\pi}{\omega}\Bigl(1-\frac{3\lambda}{4\omega^4}E+O(E^2)\Bigr).
\]
The closure rewrite follows from $\omega^2=\mu c_2$ and $\lambda=\mu c_4$. The isochrony and monotonicity statements are immediate from the sign of $\lambda$.
\end{proof}


\subsection*{3.4 Sixth extension: even hierarchy and delayed probe drift}
\addcontentsline{toc}{subsection}{3.4 Sixth extension: even hierarchy and delayed probe drift}
The quartic shell is the first operational nonlinear readout, but the reversible one-field structure already says more: if that quartic channel vanishes, the next admissible visible correction cannot appear arbitrarily. Because the exact branch potential is even, the higher hierarchy stays purely even, and the first nonzero higher closure coefficient delays the first nonzero timing defect to a correspondingly higher probe order.

\begin{theorem}[3.4 higher even hierarchy and delayed probe drift]
\label{thm:mainline-345}
Assume the hypotheses of Theorem~\ref{thm:mainline-343}. Assume moreover that on a symmetric admissible neighborhood $J$ the transported closure expansion has the form
\[
  \Gamma_{\mathrm{cl}}(\eta)
  =\Gamma_{\mathrm{cl}}(0)+\frac{c_2}{2}\eta^2+\frac{c_{2m}}{2m}\eta^{2m}+O(\eta^{2m+2}),
  \qquad m\ge 3,
\]
with
\[
  c_2>0,
  \qquad
  c_4=c_6=\cdots=c_{2m-2}=0,
  \qquad
  c_{2m}\neq 0,
\]
and that the exact visible branch potential is the transported readout of the same closure geometry up to order $2m$:
\[
  \mathcal U(\eta)=\mu\bigl(\Gamma_{\mathrm{cl}}(\eta)-\Gamma_{\mathrm{cl}}(0)\bigr)+O(\eta^{2m+2}),
  \qquad \mu>0.
\]
Then the visible branch law has the form
\[
  \mathcal U(\eta)=\frac12\omega^2\eta^2+\frac{\alpha_{2m}}{2m}\eta^{2m}+O(\eta^{2m+2}),
  \qquad
  \omega^2=\mu c_2,
  \qquad
  \alpha_{2m}=\mu c_{2m},
\]
and hence
\[
  \eta''+\omega^2\eta+\alpha_{2m}\eta^{2m-1}=O(\eta^{2m+1}).
\]
Moreover there exists a constant $\kappa_m(\omega)>0$, depending only on $m$ and $\omega$, such that the minimal period $T(E)$ of the corresponding regular small-energy periodic branch satisfies
\[
  T(E)=\frac{2\pi}{\omega}\Bigl(1-\kappa_m(\omega)\,\alpha_{2m}\,E^{m-1}+O(E^m)\Bigr)
  \qquad (E\downarrow 0).
\]
Equivalently,
\[
  T(E)=\frac{2\pi}{\omega}\Bigl(1-\kappa_m(\omega)\,\mu c_{2m}\,E^{m-1}+O(E^m)\Bigr).
\]
Consequently:
\begin{enumerate}[label=(\roman*)]
  \item the first nonlinear period correction vanishes at all lower orders $E,\dots,E^{m-2}$;
  \item the first nonzero higher timing defect is governed exactly by the first nonvanishing higher even closure coefficient $c_{2m}$;
  \item if $c_{2m}>0$, the period decreases at order $E^{m-1}$, while if $c_{2m}<0$, it increases at that order.
\end{enumerate}
The supporting normal-form reduction is recorded in Appendix~C as the sixth 3.4 extension.
\end{theorem}

\begin{proof}
Evenness of the exact branch potential was established in Theorem~\ref{thm:mainline-341}. Because the closure and visible potential agree up to order $2m$, coefficient comparison gives
\[
  \omega^2=\mu c_2,
  \qquad
  \alpha_{2m}=\mu c_{2m},
\]
and the intermediate even coefficients vanish on the visible side together with the corresponding closure coefficients. Differentiating yields the branch law
\[
  \eta''+\omega^2\eta+\alpha_{2m}\eta^{2m-1}=O(\eta^{2m+1}).
\]
A standard Poincar\'e--Lindstedt / resonant normal-form expansion for a reversible oscillator with first nonlinearity of order $2m-1$ gives a frequency law
\[
  \Omega(E)=\omega+\widetilde\kappa_m(\omega)\,\alpha_{2m}\,E^{m-1}+O(E^m),
\]
with $\widetilde\kappa_m(\omega)>0$ depending only on $m$ and $\omega$. Inverting the frequency gives the stated period expansion with $\kappa_m(\omega)>0$. Since all lower even coefficients vanish, no lower-order period correction can appear. The sign of the first nonzero timing drift is therefore exactly the sign of $\alpha_{2m}$, hence of $c_{2m}$ because $\mu>0$.
\end{proof}

\subsection*{3.4 Seventh extension: canonical even generator family}
\addcontentsline{toc}{subsection}{3.4 Seventh extension: canonical even generator family}
The delayed even hierarchy is already enough to package the visible branch law into a single local generator: if the transported closure germ is even and analytic, then the visible potential is not merely an even series term by term, but a single analytic function of $\eta^2$. This turns the 3.4 hierarchy into a canonical one-field generator family rather than a list of unrelated higher corrections.

\begin{theorem}[3.4 canonical even generator family]
\label{thm:mainline-346}
Assume the hypotheses of Theorem~\ref{thm:mainline-345}. Assume moreover that on a symmetric admissible neighborhood $J=(-R,R)$ the transported closure germ is real-analytic and even, so that
\[
  \Gamma_{\mathrm{cl}}(\eta)
  =\Gamma_{\mathrm{cl}}(0)+\sum_{m\ge 1}\frac{c_{2m}}{2m}\eta^{2m}
\]
converges for $|\eta|<R$, with $c_2>0$, and that the exact visible branch potential is the transported readout of the same germ on a possibly smaller symmetric neighborhood $J_v=(-R_v,R_v)$:
\[
  \mathcal U(\eta)=\mu\bigl(\Gamma_{\mathrm{cl}}(\eta)-\Gamma_{\mathrm{cl}}(0)\bigr),
  \qquad \mu>0,
  \qquad |\eta|<R_v\le R.
\]
Then there exists a unique real-analytic generator
\[
  \Phi(s)=\sum_{m\ge 1} a_m s^m,
  \qquad a_m=\mu\frac{c_{2m}}{2m},
  \qquad 0\le s<R_v^2,
\]
such that
\[
  \mathcal U(\eta)=\Phi(\eta^2)
  \qquad (|\eta|<R_v).
\]
Equivalently, the exact branch law has the canonical generator form
\[
  \eta'' = B(\eta) = -2\eta\,\Phi'(\eta^2),
\]
or, in expanded form,
\[
  \eta'' + \omega^2\eta + \lambda\eta^3 + \sum_{m\ge 3}\alpha_{2m}\eta^{2m-1}=0,
\]
with
\[
  \omega^2=2a_1=\mu c_2,
  \qquad
  \lambda=4a_2=\mu c_4,
  \qquad
  \alpha_{2m}=2m\,a_m=\mu c_{2m}\quad (m\ge 3).
\]
Consequently, every local nonlinear timing coefficient and every regular turning-value correction on sufficiently small admissible shells is determined by the single generator $\Phi$; in particular no second local scalar channel is opened by passing from the coefficient hierarchy to the analytic generator form.
The supporting coefficient-transport statement is recorded in Appendix~C as the seventh 3.4 extension.
\end{theorem}

\begin{proof}
Every even real-analytic germ can be written uniquely as an analytic function of $\eta^2$: grouping the convergent even series by powers of $s=\eta^2$ gives
\[
  \Gamma_{\mathrm{cl}}(\eta)-\Gamma_{\mathrm{cl}}(0)
  =\sum_{m\ge 1}\frac{c_{2m}}{2m}(\eta^2)^m.
\]
Multiplying by $\mu$ and setting
\[
  \Phi(s):=\sum_{m\ge 1}\mu\frac{c_{2m}}{2m}s^m
\]
produces a unique real-analytic generator on $[0,R_v^2)$ with $\mathcal U(\eta)=\Phi(\eta^2)$. Differentiating yields
\[
  B(\eta)=-\mathcal U'(\eta)=-2\eta\,\Phi'(\eta^2).
\]
Expanding $\Phi'(\eta^2)$ by powers of $\eta^2$ gives the stated coefficient identifications. Since Theorems~\ref{thm:mainline-342}, \ref{thm:mainline-344}, and \ref{thm:mainline-345} already identify the local shell data through the visible potential coefficients, replacing the hierarchy by the single analytic generator does not introduce a new scalar channel; it merely repackages the same transported closure data in canonical form.
\end{proof}


\subsection*{3.4 Eighth extension: generator-shell probe classifier}
\addcontentsline{toc}{subsection}{3.4 Eighth extension: generator-shell probe classifier}
The analytic generator form now lets the probe dichotomy be read directly from the shell geometry. Once the visible branch law is encoded by a single even generator $\Phi$, every sufficiently small regular shell is determined by one positive generator value and therefore carries only the transported sign ambiguity. The probe decision is then no longer an additional overlay: it is the $\mathbb Z_2$ shell classifier attached to the same generator family.

\begin{theorem}[3.4 generator-shell classifier for probe transport]
\label{thm:mainline-347}
Assume the hypotheses of Theorem~\ref{thm:mainline-346}. Assume moreover that on $[0,\rho)$ with $0<\rho<R_v^2$ one has
\[
  \Phi'(s)>0,
\]
so that the generator is strictly increasing on the small-shell sector. Then:
\begin{enumerate}[label=(\roman*)]
  \item for every regular shell energy $E\in(0,\Phi(\rho))$ there exists a unique generator radius
  \[
    s(E)\in(0,\rho)
    \qquad\text{with}\qquad
    \Phi\bigl(s(E)\bigr)=E;
  \]
  \item the corresponding visible turning pair is exactly
  \[
    \eta=\pm\sqrt{s(E)};
  \]
  \item every admissible probe transport restricted to that shell is classified by the induced sign action on the turning pair: sign-preserving transport gives direct visible convergence, while sign-reversing transport gives closure only after the unique spinorial/M\"obius lift;
  \item if no monodromy higher than the transported $\mathbb Z_2$-class opens on the shell, then this sign classifier is complete: there is no third local probe outcome on sufficiently small regular shells.
\end{enumerate}
Equivalently, on the small-shell sector the hybrid reversal probe factors through the generator-shell map
\[
  \mathfrak S_{\mathrm{rev}}
  \longrightarrow
  E
  \longmapsto
  s(E)
  \longmapsto
  \{\pm\sqrt{s(E)}\},
\]
with the final ambiguity exhausted by the transported $\mathbb Z_2$ sign action.
The supporting shell-classifier statement is recorded in Appendix~C as the eighth 3.4 extension.
\end{theorem}

\begin{proof}
Since $\Phi'(s)>0$ on $[0,\rho)$, the generator is strictly increasing there, hence invertible onto its image. This gives the unique shell parameter $s(E)$ for every small regular energy. Because $\mathcal U(\eta)=\Phi(\eta^2)$, the turning equation $\mathcal U(\eta)=E$ is equivalent to $\eta^2=s(E)$, so the visible shell consists exactly of the pair $\pm\sqrt{s(E)}$. Any admissible probe transport on that shell must therefore act on this pair by either preserving or reversing sign. The sign-preserving case closes on the same visible branch, while the sign-reversing case closes only after the already identified spinorial/M\"obius double-cover lift. If no higher monodromy is available, these two sign actions exhaust the possible local shell outcomes.
\end{proof}



\subsection*{3.4 Ninth extension: shell-parity criterion and diagnostic Julia plot}
\addcontentsline{toc}{subsection}{3.4 Ninth extension: shell-parity criterion and diagnostic Julia plot}
The generator-shell classifier can now be sharpened into a direct parity law for the local probe decision. On a fixed sufficiently small regular shell, the visible data is exhausted by the transported sign pair $\{\pm\sqrt{s(E)}\}$, so any admissible closure cycle acts by a two-element shell parity rather than by a free additional branch label. The same two-way split is also visualized by a diagnostic Julia-type render recorded in Appendix~C, now shown at the currently relevant residual scales $\gamma\in\{0,\Delta_\beta,\Delta_\theta,\Delta\mu\}$ rather than at the earlier standalone hypothesis marker $0.0009$.

\begin{theorem}[3.4 shell-parity criterion for direct closure versus spinorial lift]
\label{thm:mainline-348}
Assume the hypotheses of Theorem~\ref{thm:mainline-347}. Fix a sufficiently small regular shell of energy $E$ and let
\[
  \Sigma_E:=\{\pm\sqrt{s(E)}\}
\]
be its visible turning pair. Assume moreover that every admissible continuation cycle on that shell admits a decomposition into regular segments and finitely many transverse reflector crossings, each transverse reflector crossing implementing the sign involution
\[
  \eta\longmapsto -\eta
\]
on $\Sigma_E$. Then there exists a shell parity
\[
  \pi(E)\in\{+1,-1\}
\]
such that after one admissible closure cycle
\[
  \eta_{\mathrm{out}}=\pi(E)\,\eta_{\mathrm{in}}
\qquad
\text{for every }
\eta_{\mathrm{in}}\in\Sigma_E.
\]
More precisely:
\begin{enumerate}[label=(\roman*)]
  \item $\pi(E)=+1$ if and only if the shell closes by direct visible convergence;
  \item $\pi(E)=-1$ if and only if visible closure is reached only after the unique spinorial/M\"obius lift;
  \item if $N_E$ denotes the number of transverse reflector crossings in one admissible shell cycle, then
  \[
    \pi(E)=(-1)^{N_E}.
  \]
  Hence odd crossing parity forces sign-reversing transport, whereas even crossing parity gives direct visible closure.
\end{enumerate}
Thus the local probe outcome on a small regular shell is determined by shell parity alone once the generator radius $s(E)$ is fixed.
The corresponding diagnostic render is recorded in Appendix~C as the ninth 3.4 extension.
\end{theorem}

\begin{proof}
By Theorem~\ref{thm:mainline-347}, the shell is reduced to the visible pair $\Sigma_E=\{\pm\sqrt{s(E)}\}$ and no higher monodromy opens on the local shell sector. Therefore every admissible continuation cycle acts on $\Sigma_E$ by an element of the two-element sign group, which we denote by $\pi(E)\in\{+1,-1\}$. If $\pi(E)=+1$, then the transported representative returns to itself and the shell closes by direct visible convergence. If $\pi(E)=-1$, then the representative returns with reversed sign and visible closure occurs only after the already identified spinorial/M\"obius double-cover lift. Under the stated decomposition hypothesis, each transverse reflector crossing contributes exactly one sign involution, while regular segments preserve sign. The total action is therefore the product of $N_E$ copies of $-1$, namely $\pi(E)=(-1)^{N_E}$.
\end{proof}



\subsection*{3.4 Tenth extension: cautious detector-readout hypothesis}
\addcontentsline{toc}{subsection}{3.4 Tenth extension: cautious detector-readout hypothesis}
The diagnostic M\"obius-reflected render suggests a further hypothesis that should be stated carefully but not silently ignored. The present probe map was not tuned to detector imagery, yet it produces a visible geometry with a central transport corridor, directed opening, and boundary intensification. Within the standing phase-shift / frame-shift reading of the Companion, this motivates the following cautious interpretation: some detector-like photon images may be readout-relative shadows of the same branch-sensitive residual geometry rather than direct images of the intrinsic kernel itself.

\paragraph{Hypothesis 3.4-H (frame-shifted detector readout).}
Assume that the intrinsic kernel is observed only through an admissible shifted readout frame. Then a detector image need not reproduce the intrinsic kernel directly; it may instead display the transported boundary geometry of the residual sector. Under that reading, the central tunnel in the diagnostic render is interpreted as a transported readout channel and the bright shell frontier as a readout-sensitive boundary layer. The current claim is \emph{not} that the diagnostic render identifies a physical event, but only that the observed visual similarity is strong enough to justify a testable structural hypothesis.

\paragraph{Immediate test observables.}
The hypothesis may be probed without committing to a full detector model by monitoring diagnostic quantities such as the tunnel width $w_{\mathrm{tun}}(\gamma)$, the edge-intensification profile $I_{\mathrm{edge}}(\gamma)$, the effective opening angle $\Theta_{\mathrm{open}}(\gamma)$, and the shell-parity response $P_{\mathbb Z_2}(E)$. A stable correspondence in such observables would strengthen the readout hypothesis; failure of correspondence would simply falsify it. Thus the visual proximity is recorded here as an explicit testable hypothesis rather than as an identification claim.

\subsection*{3.4 Eleventh extension: tunnel observables and symmetry-constrained readout metrics}
\addcontentsline{toc}{subsection}{3.4 Eleventh extension: tunnel observables and symmetry-constrained readout metrics}
The detector-readout hypothesis becomes mathematically useful only once the visual geometry is reduced to explicit observables. The diagnostic render therefore should be read through a small family of shell metrics attached to the same generator/parity structure rather than through unaudited visual impression alone. In particular, the central transport corridor, the shell edge brightening, and the opening of the visible branch front can be organized into a minimal test package.

\begin{proposition}[3.4 tunnel-observable package and parity-compatible readout]
\label{prop:mainline-3410}
Fix a diagnostic probe family $\mathcal J_\gamma$ on a symmetric rendering window $W\subset\mathbb C$ together with a nonnegative stability-density image
\[
  H_\gamma:W\to\mathbb R_{\ge 0}
\]
obtained from the corresponding escape-time or residence-time readout. Assume that the rendering window and the probe normalization are chosen so that the transported shell classifier of Theorem~\ref{thm:mainline-348} is visible on a distinguished symmetry axis $x=0$. Then the following four scalar observables are canonically available on the small-shell diagnostic sector:
\begin{enumerate}[label=(\roman*)]
  \item the \emph{tunnel width}
  \[
    w_{\mathrm{tun}}(\gamma;\tau)
    :=\operatorname{meas}\bigl\{y:(0,y)\in W,\ H_\gamma(0,y)\le \tau\bigr\},
  \]
  defined for any admissible intensity threshold $\tau>0$;
  \item the \emph{edge-intensification functional}
  \[
    I_{\mathrm{edge}}(\gamma)
    :=\int_{\partial S_\gamma} H_\gamma\,d\sigma,
  \]
  where $S_\gamma$ is the first regular shell selected by the generator-radius prescription of Theorem~\ref{thm:mainline-347};
  \item the \emph{opening angle}
  \[
    \Theta_{\mathrm{open}}(\gamma),
  \]
  defined as the minimal cone angle based at the central channel that contains the first visible shell frontier on the upper half-window;
  \item the already available shell-parity response
  \[
    P_{\mathbb Z_2}(E)=\pi(E)\in\{+1,-1\}.
  \]
\end{enumerate}
If, in addition, the diagnostic family is normalized so that the transported sign action is the only local monodromy on sufficiently small shells, then the quadruple
\[
  \mathfrak O_{\mathrm{diag}}(\gamma)
  :=\bigl(w_{\mathrm{tun}}(\gamma;\tau),\ I_{\mathrm{edge}}(\gamma),\ \Theta_{\mathrm{open}}(\gamma),\ P_{\mathbb Z_2}(E)\bigr)
\]
contains no independent local scalar beyond the generator-shell data: each component is a readout-functional of the same shell radius $s(E)$ and parity class $\pi(E)$, possibly after the external choice of the plotting threshold $\tau$.
\end{proposition}

\begin{proof}
Theorem~\ref{thm:mainline-347} reduces every sufficiently small regular shell to the generator radius $s(E)$ and the transported sign pair $\{\pm\sqrt{s(E)}\}$, while Theorem~\ref{thm:mainline-348} shows that the only remaining local branch ambiguity is the parity $\pi(E)\in\{+1,-1\}$. Once a concrete diagnostic density $H_\gamma$ is fixed, the corridor width, shell-edge intensity, and opening angle are ordinary scalar functionals of the rendered shell geometry. Since the shell geometry itself is determined by $s(E)$ together with the parity action, these observables do not open a new local scalar channel; they only read the same generator/parity data through different diagnostics. The threshold $\tau$ is external plotting metadata and therefore does not count as an intrinsic structural scalar.
\end{proof}

\paragraph{Operational consequence.}
The proposition does not claim that the four observables already agree with detector data; it only shows that the detector-readout hypothesis can be made auditably quantitative. In particular, the central tunnel noted in the current diagnostic render is not merely a visual remark: after thresholding it becomes the measurable quantity $w_{\mathrm{tun}}(\gamma;\tau)$, while the shell brightening and opening become $I_{\mathrm{edge}}(\gamma)$ and $\Theta_{\mathrm{open}}(\gamma)$. This is the first point at which the image layer becomes a genuine test interface rather than a suggestive illustration.


\subsection*{3.4 Twelfth extension: calibrated inverse readout and shell-consistency defect}
\addcontentsline{toc}{subsection}{3.4 Twelfth extension: calibrated inverse readout and shell-consistency defect}
The quantitative tunnel package becomes substantially stronger once one asks not only how shell data produces visible metrics, but also how far such metrics may be inverted. The natural minimal requirement is a local calibration law: tunnel width should move monotonically with the generator shell radius on a fixed threshold protocol. Under that assumption, the image layer becomes not merely comparable but locally readable.

\begin{theorem}[3.4 calibrated inverse readout and shell consistency]
\label{thm:mainline-3411}
Fix a threshold level $\tau>0$ and assume the diagnostic family of Proposition~\ref{prop:mainline-3410} is restricted to a sufficiently small regular shell sector on which:
\begin{enumerate}[label=(\roman*)]
  \item the tunnel-width functional factors through the generator shell radius,
  \[
    w_{\mathrm{tun}}(\gamma;\tau)=W_\tau(s(E)),
  \]
  for a $C^1$ calibration map $W_\tau$;
  \item $W_\tau$ is strictly monotone on the relevant shell interval;
  \item the parity response records the transported sign class,
  \[
    P_{\mathbb Z_2}(E)=\pi(E)\in\{+1,-1\}.
  \]
\end{enumerate}
Then the pair
\[
  \bigl(w_{\mathrm{tun}}(\gamma;\tau),P_{\mathbb Z_2}(E)\bigr)
\]
locally determines the generator/parity data
\[
  \bigl(s(E),\pi(E)\bigr)
  =\bigl(W_\tau^{-1}(w_{\mathrm{tun}}(\gamma;\tau)),P_{\mathbb Z_2}(E)\bigr).
\]
Consequently, the remaining observables of Proposition~\ref{prop:mainline-3410} admit local consistency representations
\[
  I_{\mathrm{edge}}(\gamma)=\widehat{\mathcal I}_\tau\bigl(w_{\mathrm{tun}}(\gamma;\tau),P_{\mathbb Z_2}(E)\bigr),
\]
\[
  \Theta_{\mathrm{open}}(\gamma)=\widehat{\vartheta}_\tau\bigl(w_{\mathrm{tun}}(\gamma;\tau),P_{\mathbb Z_2}(E)\bigr)
\]
for suitable local readout functions $\widehat{\mathcal I}_\tau$ and $\widehat{\vartheta}_\tau$. In particular, no new local structural scalar opens at the inverse-readout level.
\end{theorem}

\begin{proof}
Strict monotonicity makes $W_\tau$ locally invertible, so the tunnel width recovers the shell radius $s(E)$ uniquely on the selected shell interval. The parity response already equals the transported sign class by assumption. The shell geometry on the small diagnostic sector is determined by $(s(E),\pi(E))$ through Theorem~\ref{thm:mainline-348} and Proposition~\ref{prop:mainline-3410}; therefore every additional shell metric constructed from the same thresholded rendering protocol factors through the recovered pair. This yields the two stated consistency maps and shows that inverse readout does not create a new intrinsic scalar channel.
\end{proof}

\begin{corollary}[3.4 shell-consistency defect]
\label{cor:mainline-3411-defect}
Under the hypotheses of Theorem~\ref{thm:mainline-3411}, every calibrated diagnostic image carries a local shell-consistency defect
\[
  \Delta_{\mathrm{diag}}(\gamma;\tau)
  :=\Bigl|I_{\mathrm{edge}}(\gamma)-\widehat{\mathcal I}_\tau\bigl(w_{\mathrm{tun}}(\gamma;\tau),P_{\mathbb Z_2}(E)\bigr)\Bigr|
   +\Bigl|\Theta_{\mathrm{open}}(\gamma)-\widehat{\vartheta}_\tau\bigl(w_{\mathrm{tun}}(\gamma;\tau),P_{\mathbb Z_2}(E)\bigr)\Bigr|.
\]
Hence $\Delta_{\mathrm{diag}}$ vanishes on an idealized perfectly calibrated render and measures the failure of the observed image to close back to the shell/parity model.
\end{corollary}

\begin{proof}
This is immediate from the consistency representations of Theorem~\ref{thm:mainline-3411}. The quantity $\Delta_{\mathrm{diag}}$ is simply the residual mismatch between the measured edge/opening readouts and the values predicted from the recovered shell/parity data.
\end{proof}

\paragraph{Interpretive consequence.}
The image layer now supports a minimal inverse logic. After choosing a threshold protocol, the tunnel width and the shell parity are enough to reconstruct the local shell radius and to predict the remaining edge/opening observables. This does not yet prove the detector-readout hypothesis, but it upgrades the diagnostic render from a suggestive comparison picture to a partially invertible readout device.



\subsection*{3.4 Thirteenth extension: inverse identifiability and stability margin}
\addcontentsline{toc}{subsection}{3.4 Thirteenth extension: inverse identifiability and stability margin}
The inverse-readout theorem becomes substantially more useful once one knows when the reconstruction is not merely available but also stable. The right local criterion is that the calibration law should have nonvanishing derivative on the shell under consideration. Then the tunnel width is not just invertible in principle: it controls the shell radius with an explicit local conditioning constant.

\begin{theorem}[3.4 local inverse identifiability and stability]
\label{thm:mainline-3412}
Assume the hypotheses of Theorem~\ref{thm:mainline-3411}. Fix a regular shell $E_0$ with generator radius
\[
  s_0:=s(E_0),
  \qquad
  \pi_0:=\pi(E_0),
  \qquad
  w_0:=W_\tau(s_0).
\]
Assume moreover that $W_\tau\in C^1$ and
\[
  W_\tau'(s_0)\neq 0.
\]
Then there exist neighborhoods $U_{s_0}$ of $s_0$ and $V_{w_0}$ of $w_0$, together with constants $m_0,M_0>0$, such that
\[
  m_0\,|s_1-s_2|
  \le
  |W_\tau(s_1)-W_\tau(s_2)|
  \le
  M_0\,|s_1-s_2|
\]
for all $s_1,s_2\in U_{s_0}$. Consequently:
\begin{enumerate}[label=(\roman*)]
  \item on the fixed parity branch $\pi=\pi_0$, the inverse readout is locally unique;
  \item for every calibrated pair $(\widetilde w,\widetilde\pi)\in V_{w_0}\times\{\pi_0\}$, the reconstructed shell radius
  \[
    \widetilde s:=W_\tau^{-1}(\widetilde w)
  \]
  satisfies the stability bound
  \[
    |\widetilde s-s_0|
    \le
    \frac{1}{m_0}\,|\widetilde w-w_0|;
  \]
  \item if the consistency maps $\widehat{\mathcal I}_\tau$ and $\widehat{\vartheta}_\tau$ from Theorem~\ref{thm:mainline-3411} are locally Lipschitz on $V_{w_0}\times\{\pi_0\}$ with constants $L_I$ and $L_\Theta$, then
  \[
    \bigl|\widehat{\mathcal I}_\tau(\widetilde w,\widetilde\pi)-\widehat{\mathcal I}_\tau(w_0,\pi_0)\bigr|
    \le L_I\,|\widetilde w-w_0|,
  \]
  \[
    \bigl|\widehat{\vartheta}_\tau(\widetilde w,\widetilde\pi)-\widehat{\vartheta}_\tau(w_0,\pi_0)\bigr|
    \le L_\Theta\,|\widetilde w-w_0|.
  \]
\end{enumerate}
Hence the calibrated inverse readout is locally bi-Lipschitz on each parity branch and is destabilized only at parity changes or at points where $W_\tau'$ vanishes.
\end{theorem}

\begin{proof}
Because $W_\tau'(s_0)\neq 0$, continuity of $W_\tau'$ gives a neighborhood $U_{s_0}$ on which $|W_\tau'|$ is bounded above and below by positive constants. The mean-value theorem then yields the two-sided estimate with suitable $m_0,M_0>0$. Local injectivity and the bound for $|\widetilde s-s_0|$ follow immediately. The Lipschitz estimates for the reconstructed edge/opening readouts are obtained by composing the local inverse branch with the Lipschitz consistency maps. A local degeneracy can therefore occur only when the parity branch changes or when the calibration derivative vanishes.
\end{proof}

\begin{corollary}[3.4 inverse condition number and stability margin]
\label{cor:mainline-3412-condition}
Under the hypotheses of Theorem~\ref{thm:mainline-3412}, define the local inverse stability margin and condition number by
\[
  \mathfrak m_\tau(s_0):=|W_\tau'(s_0)|,
  \qquad
  \kappa_\tau(s_0):=\frac{1}{|W_\tau'(s_0)|}.
\]
Then larger $\mathfrak m_\tau(s_0)$ (equivalently smaller $\kappa_\tau(s_0)$) means better-conditioned shell reconstruction from tunnel width, while parity-switch points and critical points of $W_\tau$ are the only local obstructions to stable inversion.
\end{corollary}

\begin{proof}
This is just a restatement of the inverse estimate from Theorem~\ref{thm:mainline-3412}. The derivative $|W_\tau'(s_0)|$ is the local amplification factor from shell-radius error to tunnel-width error, so its reciprocal is the condition number for the inverse problem. The final sentence records the two ways in which the theorem can fail.
\end{proof}

\paragraph{Interpretive consequence.}
The image layer now comes with a local reliability scale. Small tunnel-width perturbations produce controlled shell-radius perturbations whenever the calibration slope stays away from zero, and the parity sector remains the only discrete ambiguity. In that sense the diagnostic plot is no longer just partially invertible; it is partially invertible with an explicit local stability regime.

\subsection*{3.4 Fourteenth extension: shell atlas and critical-shell stratification}
\addcontentsline{toc}{subsection}{3.4 Fourteenth extension: shell atlas and critical-shell stratification}
Local inverse stability becomes structurally much clearer once the admissible shell line is decomposed into regular inverse charts separated only by genuine degeneracy loci. The natural global objects are therefore the critical-shell set of the calibration law and the parity-switch set of the transported probe class. Away from these two obstructions, the inverse readout does not merely exist pointwise: it assembles into a shell atlas with uniform local control on compact subintervals.

\begin{theorem}[3.4 shell atlas and critical-shell stratification]
\label{thm:mainline-3413}
Assume the hypotheses of Theorem~\ref{thm:mainline-3412}. Let $I\subset \mathbb R_{\ge 0}$ be a shell interval in the generator coordinate $s$, and define the critical-shell and parity-switch sets by
\[
\begin{aligned}
\Sigma_{\mathrm{crit}}&:=\{s\in I: W_\tau'(s)=0\},\\
\Sigma_{\mathrm{par}}&:=\{s\in I:\text{the transported parity class changes at }s\}.
\end{aligned}
\]
Set
\[
  \Sigma:=\Sigma_{\mathrm{crit}}\cup\Sigma_{\mathrm{par}}.
\]
Then every connected component $J\subset I\setminus \Sigma$ carries a calibrated shell chart with the following properties:
\begin{enumerate}[label=(\roman*)]
  \item the parity class is constant on $J$;
  \item the calibration law $W_\tau|_J$ is strictly monotone, hence injective;
  \item for every compact subinterval $K\Subset J$ there exist constants $m_K,M_K>0$ such that
  \[
    m_K|s_1-s_2|
    \le |W_\tau(s_1)-W_\tau(s_2)|
    \le M_K|s_1-s_2|
  \]
  for all $s_1,s_2\in K$;
  \item on each such $K$, the inverse readout extends to a stable branch chart
  \[
    (w_{\mathrm{tun}},P_{\mathbb Z_2})
    \longmapsto
    (s,\pi)
  \]
  with Lipschitz inverse constant bounded by $1/m_K$.
\end{enumerate}
Consequently the calibrated inverse problem on $I$ is stratified by the closed obstruction set $\Sigma$: away from $\Sigma$ one has a finite or countable atlas of stable parity-preserving shell charts, and every local inverse degeneracy is forced by either a calibration critical shell or a parity-switch shell.
\end{theorem}

\begin{proof}
On every connected component $J\subset I\setminus \Sigma$, the parity is locally constant by definition of $\Sigma_{\mathrm{par}}$. Since $W_\tau'$ never vanishes on $J$, continuity of $W_\tau'$ implies that $W_\tau'$ cannot change sign on $J$; hence $W_\tau|_J$ is strictly monotone and therefore injective. For any compact $K\Subset J$, continuity of $W_\tau'$ yields positive lower and finite upper bounds on $|W_\tau'|$, and the mean-value theorem gives the two-sided estimate with constants $m_K,M_K$. The inverse-chart statement is then exactly the local inverse-stability estimate of Theorem~\ref{thm:mainline-3412}, now uniform on $K$. Since every failure of this argument occurs only when parity changes or $W_\tau'$ vanishes, the obstruction set is precisely $\Sigma$.
\end{proof}

\begin{corollary}[3.4 chartwise diagnostic regime]
\label{cor:mainline-3413-chartwise}
Under the hypotheses of Theorem~\ref{thm:mainline-3413}, the shell-consistency defect
\[
  \Delta_{\mathrm{diag}}(\gamma;\tau)
\]
becomes a chartwise diagnostic observable on every compact $K\Subset J\subset I\setminus\Sigma$. In particular, a persistent rise of $\Delta_{\mathrm{diag}}$ along a stable chart is evidence for model mismatch inside a fixed shell/parity regime, whereas loss of invertibility itself can occur only by approaching $\Sigma_{\mathrm{crit}}$ or $\Sigma_{\mathrm{par}}$.
\end{corollary}

\begin{proof}
Theorem~\ref{thm:mainline-3413} gives a stable inverse chart on every compact $K\Subset J$. Inside such a chart, the reconstructed shell/parity data vary continuously with the measured tunnel width, so the shell-consistency defect is a well-defined chart observable. A blow-up of the defect while the inverse chart remains valid indicates failure of the model closure rather than loss of invertibility. The only mechanisms that can destroy the chart itself are exactly the two obstruction loci already identified in Theorem~\ref{thm:mainline-3413}.
\end{proof}

\paragraph{Interpretive consequence.}
The inverse-readout layer now has a global local-to-global organization. Stable reconstruction is not a scattered pointwise accident, but comes in shell charts separated by sharply identified obstruction shells. This makes the diagnostic side much cleaner: one can distinguish chart-internal model mismatch from genuine loss of inverse control.

\begin{remark}[OP3 reading of shell diagnostics]
In the present release, shell-consistency defects in a stable chart are read through the OP3 residual/shell interface and the visible/structural mismatch separation tool. A rise of $\Delta_{\mathrm{diag}}$ inside a valid chart is first an OP3-admissibility or model-fit issue inside a fixed shell/parity regime, not immediate evidence against the closed OP~1--OP~5 source-return pair.
\end{remark}

\subsection*{3.4 Fifteenth extension: obstruction-shell transition law}
\addcontentsline{toc}{subsection}{3.4 Fifteenth extension: obstruction-shell transition law}
Once the inverse readout is organized by stable shell charts, the next structural question is how the data reorganizes when a trajectory reaches an obstruction shell. There are only two primitive failure modes in the present atlas: calibration-fold shells, where the tunnel-width derivative vanishes, and parity shells, where the transported $\mathbb Z_2$ class flips. The point of the next statement is that these are not arbitrary singularities; they carry a very small transition algebra.

\begin{theorem}[3.4 obstruction-shell transition law]
\label{thm:mainline-3414}
Let $s_\ast\in \Sigma=\Sigma_{\mathrm{crit}}\cup\Sigma_{\mathrm{par}}$ be an isolated obstruction shell in the sense of Theorem~\ref{thm:mainline-3413}.
\begin{enumerate}[label=(\roman*)]
  \item If $s_\ast\in \Sigma_{\mathrm{crit}}\setminus\Sigma_{\mathrm{par}}$ and
  \[
    W_\tau'(s_\ast)=0,
    \qquad
    W_\tau''(s_\ast)\neq 0,
  \]
  then the parity class is constant near $s_\ast$, and there exists a punctured neighborhood on which the calibrated tunnel law has exactly two local inverse sheets joined at $s_\ast$. More precisely, after a local reparametrization one has the fold normal form
  \[
    w_{\mathrm{tun}}-w_\ast = \pm u^2 + O(u^3),
    \qquad
    u=s-s_\ast,
  \]
  so equal tunnel width near $w_\ast$ corresponds to the two shell branches exchanged by the local involution
  \[
    \sigma_\ast(s)\neq s,
    \qquad
    W_\tau(\sigma_\ast(s))=W_\tau(s),
    \qquad
    \sigma_\ast^2=\mathrm{id}.
  \]

  \item If $s_\ast\in \Sigma_{\mathrm{par}}\setminus\Sigma_{\mathrm{crit}}$, then $W_\tau'(s_\ast)\neq 0$, hence the shell radius continues uniquely through $s_\ast$, but the transported parity flips:
  \[
    \pi_+(s_\ast)=-\pi_-(s_\ast).
  \]
  Thus the obstruction is purely discrete: the shell coordinate is continuous, while the visible branch label changes sign.

  \item If $s_\ast\in \Sigma_{\mathrm{crit}}\cap\Sigma_{\mathrm{par}}$, then the local transition is the composition of the two previous mechanisms: a fold exchange of shell sheets together with a parity flip. In particular, the local transition algebra is generated by
  \[
    T_\ast^{\mathrm{fold}}:(s,\pi)\mapsto (\sigma_\ast(s),\pi),
    \qquad
    T_\ast^{\mathrm{par}}:(s,\pi)\mapsto (s,-\pi),
  \]
  and every admissible local transition is one of
  \[
    \mathrm{id},
    \quad
    T_\ast^{\mathrm{fold}},
    \quad
    T_\ast^{\mathrm{par}},
    \quad
    T_\ast^{\mathrm{par}}\circ T_\ast^{\mathrm{fold}}.
  \]
\end{enumerate}
Hence obstruction shells do not open a new uncontrolled singular sector: they only create fold exchange, parity jump, or their composition.
\end{theorem}

\begin{proof}
For (i), the assumptions say that $s_\ast$ is a nondegenerate critical point of the scalar calibration law. The one-dimensional Morse lemma (equivalently a standard fold normal form for a smooth map with first nonvanishing second derivative) gives a local coordinate $u$ in which $W_\tau-W_\tau(s_\ast)$ is quadratic up to higher order, hence the inverse problem has exactly two sheets on one side of the critical value and none on the other. Because $s_\ast\notin \Sigma_{\mathrm{par}}$, parity is constant through this local sector, so only the shell sheet changes. The local deck involution exchanging the two branches is exactly $\sigma_\ast$.

For (ii), $s_\ast\notin \Sigma_{\mathrm{crit}}$ implies $W_\tau'(s_\ast)\neq 0$, so the inverse function theorem gives a unique continuation of the shell coordinate across $s_\ast$. Since $s_\ast\in \Sigma_{\mathrm{par}}$, the only local change is the transported parity flip.

Statement (iii) is immediate by combining the two previous local descriptions. The fold exchange leaves parity unchanged, the parity jump leaves the shell coordinate unchanged, and the two operations commute because they act on different coordinates. Therefore the local transition algebra is the Klein four-set generated by these two involutions.
\end{proof}

\begin{corollary}[3.4 obstruction monodromy principle]
\label{cor:mainline-3414}
Under the hypotheses of Theorem~\ref{thm:mainline-3414}, every closed local continuation of admissible shell data around isolated obstruction shells produces only a $\mathbb Z_2$-generated monodromy on the inverse data $(s,\pi)$. In particular, direct visible return corresponds to trivial monodromy, while the spinorial/M\"obius case is exactly the nontrivial parity component of this local obstruction monodromy.
\end{corollary}

\begin{proof}
By Theorem~\ref{thm:mainline-3414}, the only local generators are the fold involution and the parity involution. Any closed continuation is therefore represented by a finite product of commuting involutions, so only a $\mathbb Z_2$-generated monodromy can occur. The visible return is trivial exactly when the parity component is trivial; otherwise the continuation returns only after the nontrivial parity lift, which is precisely the local spinorial/M\"obius case.
\end{proof}

\paragraph{Interpretive consequence.}
The shell atlas now comes with transition rules rather than only forbidden loci. Critical shells create fold exchange, parity shells create sign change, and the mixed case composes the two. This is the first point at which the inverse image layer acquires a genuine obstruction algebra instead of a mere list of bad points.


\subsection*{3.4 Closing synthesis: generator, inverse atlas, and parity holonomy}
\addcontentsline{toc}{subsection}{3.4 Closing synthesis: generator, inverse atlas, and parity holonomy}
The 3.4 line can now be closed without adding a new mechanism. The previous extensions already determine a single visible package: an even analytic generator for the residual branch law, a shell atlas for calibrated inverse readout, and a parity obstruction that controls the only remaining global ambiguity. The point of the present synthesis is therefore not to enlarge the theory, but to compress the accumulated 3.4 structure into one handoff theorem.

\begin{theorem}[3.4 closing synthesis package]
\label{thm:mainline-3420}
Assume the hypotheses of the stabilized 3.4 line on a connected admissible shell base $B$. Then the visible residual readout is controlled by the following data and no further independent local scalar is opened by the 3.4 construction:
\begin{enumerate}[label=(\roman*), leftmargin=2em]
  \item there exists an even analytic generator $\Phi$ such that the visible branch law is
  \[
    \eta'' = -2\eta\,\Phi'(\eta^2),
    \qquad
    \mathcal U(\eta)=\Phi(\eta^2),
  \]
  and every sufficiently small regular shell is determined by the generator radius
  \[
    s(E)=\Phi^{-1}(E)
  \]
  together with the transported parity label $\pi\in\mathbb Z_2$;

  \item away from the obstruction set $\Sigma=\Sigma_{\mathrm{crit}}\cup\Sigma_{\mathrm{par}}$, the calibrated inverse readout
  \[
    (w_{\mathrm{tun}},P_{\mathbb Z_2})\longmapsto (s,\pi)
  \]
  is chartwise stable, and the remaining observables $I_{\mathrm{edge}}$ and $\Theta_{\mathrm{open}}$ are shell-consistency readouts of the same data;

  \item the global ambiguity of the inverse shell atlas is exhausted by the parity local system, equivalently by the cohomological class
  \[
    \varepsilon_{\mathrm{par}}\in H^1(B,\mathbb Z_2),
  \]
  whose loop readout is the parity holonomy character
  \[
    \chi_{\mathrm{par}}:\pi_1(B,b_0)\to\mathbb Z_2;
  \]

  \item the empirical estimator $\widehat\chi_{\mathrm{par}}$ obtained from parity-changing obstruction crossings recovers this same global class along admissible closed shell traces.
\end{enumerate}
Consequently, direct visible closure is globally available exactly in the trivial parity class, while every nontrivial case becomes globally single-valued on the canonical parity double cover.
\end{theorem}

\begin{proof}
Item (i) is exactly the content of the generator reduction established in Theorems~\ref{thm:mainline-346} and~\ref{thm:mainline-347}. Item (ii) combines the calibrated inverse readout theorem, the inverse identifiability/stability theorem, and the shell-atlas stratification from Theorems~\ref{thm:mainline-3411}, \ref{thm:mainline-3412}, and~\ref{thm:mainline-3413}. Item (iii) is the cohomological and double-cover compression proved in Theorems~\ref{thm:mainline-3415}--\ref{thm:mainline-3418}. Item (iv) is precisely the empirical parity-holonomy estimator of Theorem~\ref{thm:mainline-3419}. No new local scalar is introduced in this synthesis step; it only packages the already established 3.4 data into a single visible readout layer.
\end{proof}

\begin{corollary}[3.4 handoff package for the next stage]
\label{cor:mainline-3420}
At the end of the 3.4 line, the reduced visible input needed for the next structural stage consists of the generator germ $\Phi$, the shell-consistency defect $\Delta_{\mathrm{diag}}$, and the parity obstruction package $(\varepsilon_{\mathrm{par}},\chi_{\mathrm{par}},\widehat\chi_{\mathrm{par}})$. In particular, the transition to the next stage does not require a new visible branch variable beyond $\eta$.
\end{corollary}

\begin{proof}
Theorem~\ref{thm:mainline-3420} already identifies the complete visible 3.4 package. The generator controls the local residual law and regular shell data, the defect term measures chartwise mismatch of the diagnostic readout, and the parity package records the only remaining global ambiguity. Therefore these data form the minimal handoff set to the next stage.
\end{proof}

\paragraph{3.4 closing remark.}
The 3.4 line has therefore reached a natural stopping point. It began by globalizing the local residual law and ends with a parity-resolved inverse readout atlas whose only global obstruction is a single class in $H^1(B,\mathbb Z_2)$. This is the right place to stop extending 3.4 and hand the package forward, rather than continuing to add more local refinements.



\subsection*{3.5 First extension: admissible residual record and reconstruction normal form}
\addcontentsline{toc}{subsection}{3.5 First extension: admissible residual record and reconstruction normal form}
The closed 3.4 line already shows that the visible branch layer is controlled by a small package rather than by an open-ended family of new observables. The correct first move of 3.5 is therefore not to enlarge the visible theory, but to treat this package as a single admissible record. This gives a compact interface for later simplification work while remaining entirely inside the current mathematical language.

\begin{definition}[Admissible residual record]
\label{def:mainline-3500}
On a connected admissible shell base $B$, define the visible residual record
\[
  \mathfrak R(B)
  :=
  \bigl(\Phi,\Delta_{\mathrm{diag}},\varepsilon_{\mathrm{par}},\chi_{\mathrm{par}},\widehat\chi_{\mathrm{par}}\bigr),
\]
where $\Phi$ is the even analytic generator of the local visible branch law, $\Delta_{\mathrm{diag}}$ is the shell-consistency defect, $\varepsilon_{\mathrm{par}}\in H^1(B,\mathbb Z_2)$ is the parity obstruction class, $\chi_{\mathrm{par}}$ is its parity-holonomy character, and $\widehat\chi_{\mathrm{par}}$ is the empirical loop estimator.
\end{definition}

\begin{theorem}[3.5 reconstruction normal form]
\label{thm:mainline-3500}
Assume the hypotheses of Theorem~\ref{thm:mainline-3420}. Then every visible statement in the stabilized generator-atlas-holonomy sector factors through the admissible residual record $\mathfrak R(B)$ together with the chartwise shell coordinate $s$ and parity label $\pi$, in the following sense:
\begin{enumerate}[label=(\roman*), leftmargin=2em]
  \item on each regular shell chart away from the obstruction set, the local branch law, tunnel observables, and shell-consistency diagnostics are determined by $(\Phi,s,\pi,\Delta_{\mathrm{diag}})$;
  \item globally, the only remaining ambiguity of visible continuation is determined by the parity obstruction package $(\varepsilon_{\mathrm{par}},\chi_{\mathrm{par}},\widehat\chi_{\mathrm{par}})$;
  \item no further independent visible branch datum beyond $\eta$ is required to reconstruct the 3.4 readout sector.
\end{enumerate}
Equivalently, the stabilized visible branch layer of the Companion is already in reconstruction normal form once it is expressed through $\mathfrak R(B)$.
\end{theorem}

\begin{proof}
Theorem~\ref{thm:mainline-3420} identifies the complete 3.4 handoff package and states that no further independent local scalar is opened by the 3.4 construction. Its item~(i) reduces the local visible residual dynamics to the generator $\Phi$, while item~(ii) reduces the calibrated inverse shell readout to chartwise shell data $(s,\pi)$ together with the shell-consistency defect. Items~(iii) and~(iv) show that the only remaining global ambiguity is the parity obstruction package and that its operational loop readout is completely captured by $\chi_{\mathrm{par}}$ and $\widehat\chi_{\mathrm{par}}$. Therefore every visible statement proved in the stabilized 3.4 sector factors through the single record $\mathfrak R(B)$ together with chartwise shell coordinates, and no new visible branch datum beyond $\eta$ is needed.
\end{proof}

\begin{corollary}[3.5 starting principle: record-level compression]
\label{cor:mainline-3500}
For the first simplification stage of v\docversion, proofs and comparisons may treat $\mathfrak R(B)$ as the complete visible residual interface. In particular, any later compact formal layer that is introduced for proof economy must remain faithful to the reconstruction normal form determined by $\mathfrak R(B)$.
\end{corollary}

\begin{proof}
This is a direct restatement of Theorem~\ref{thm:mainline-3500}. Once the visible branch layer has been compressed to $\mathfrak R(B)$, any further simplification may only reorganize or repackage these same data, not enlarge the visible branch interface by adding new primitive observables.
\end{proof}

\paragraph{3.5 opening remark.}
The transition from 3.4 to 3.5 is therefore a change of stance rather than a new visible mechanism. Instead of extending the branch picture by more local refinements, the document now treats the residual generator, the inverse shell atlas, and the parity holonomy package as one admissible record. This is the first point at which later proof-economy work can aim at genuine compression instead of further accumulation.

\subsection*{3.5 Second extension: record morphisms and faithful translation principle}
\addcontentsline{toc}{subsection}{3.5 Second extension: record morphisms and faithful translation principle}
Once the visible residual sector has been compressed to an admissible record, the next question is no longer how to add more local data, but how to compare two residual presentations without reopening the visible layer. The correct next step is therefore to formulate a record-level notion of sameness. This does not yet introduce a new formal language; it only states which data are allowed to move and which must remain invariant under later compression.

\begin{definition}[Admissible residual-record morphism]
\label{def:mainline-3501}
Let $B_1$ and $B_2$ be connected admissible shell bases with residual records
\[
  \mathfrak R(B_i)=\bigl(\Phi_i,\Delta^{(i)}_{\mathrm{diag}},\varepsilon^{(i)}_{\mathrm{par}},\chi^{(i)}_{\mathrm{par}},\widehat\chi^{(i)}_{\mathrm{par}}\bigr),
  \qquad i=1,2.
\]
An \emph{admissible residual-record morphism} $F:B_1\to B_2$ is an admissible shell-atlas transport such that, on every regular shell chart away from the obstruction set, the following conditions hold:
\begin{enumerate}[label=(\roman*), leftmargin=2em]
  \item the generator germ is preserved under transport, equivalently the visible branch law is carried by the same transported generator data;
  \item the shell-consistency defect is preserved up to chartwise pullback,
  \[
    \Delta^{(1)}_{\mathrm{diag}} = F^*\Delta^{(2)}_{\mathrm{diag}};
  \]
  \item the parity obstruction package is preserved,
  \[
    \varepsilon^{(1)}_{\mathrm{par}} = F^*\varepsilon^{(2)}_{\mathrm{par}},
    \qquad
    \chi^{(1)}_{\mathrm{par}} = \chi^{(2)}_{\mathrm{par}}\circ F_*,
    \qquad
    \widehat\chi^{(1)}_{\mathrm{par}} = \widehat\chi^{(2)}_{\mathrm{par}}\circ F_*.
  \]
\end{enumerate}
If, moreover, $F$ is invertible in the same class, then $F$ is called a \emph{residual-record isomorphism}.
\end{definition}

\begin{theorem}[3.5 faithful translation principle]
\label{thm:mainline-3501}
Assume the hypotheses of Theorem~\ref{thm:mainline-3500}. Let $F:B_1\to B_2$ be a residual-record isomorphism in the sense of Definition~\ref{def:mainline-3501}. Then every visible statement in the stabilized generator-atlas-holonomy sector is invariant under translation along $F$. More precisely:
\begin{enumerate}[label=(\roman*), leftmargin=2em]
  \item local branch laws, shell reconstructions, and tunnel observables translate chartwise without opening a new visible scalar channel;
  \item the shell-consistency defect translates by pullback, so chartwise mismatch diagnostics are record-invariant;
  \item the parity obstruction class, its holonomy character, and the empirical loop estimator are preserved under transport of admissible loops;
  \item consequently, any theorem proved solely from the stabilized 3.4 package depends only on the residual-record isomorphism class of $\mathfrak R(B)$.
\end{enumerate}
Equivalently, the visible generator-atlas-holonomy sector is faithful to record level: later proof compression may reorganize it, but may not change its record-level content.
\end{theorem}

\begin{proof}
Theorem~\ref{thm:mainline-3500} already states that every visible statement in the stabilized sector factors through the record $\mathfrak R(B)$ together with chartwise shell coordinate and parity label. By Definition~\ref{def:mainline-3501}, a residual-record isomorphism preserves exactly the five pieces of data entering this factorization: the transported generator germ, the shell-consistency defect, the cohomological parity class, the holonomy character, and the empirical estimator. Hence local branch formulas, inverse-shell reconstructions, and calibrated tunnel observables translate chartwise without opening a new visible scalar. Since the parity package is preserved under loop transport, the global ambiguity sector is likewise invariant. Therefore any theorem whose hypotheses and conclusions are expressed entirely in the stabilized generator-atlas-holonomy sector depends only on the record-level isomorphism class and is faithfully transportable along $F$.
\end{proof}

\begin{corollary}[3.5 translation corollary: record-level theorem transport]
\label{cor:mainline-3501}
Let $\mathcal S_1$ and $\mathcal S_2$ be two admissible shell presentations related by a residual-record isomorphism. Then any proof in the stabilized visible sector may be transported from $\mathcal S_1$ to $\mathcal S_2$ without introducing a new primitive observable, provided the proof uses only record-level data and chartwise shell coordinates.
\end{corollary}

\begin{proof}
This is the direct content of Theorem~\ref{thm:mainline-3501}. Once the stabilized sector has been reduced to record level, proof transport between isomorphic shell presentations is faithful and does not enlarge the visible interface.
\end{proof}

\paragraph{3.5.1 working remark.}
The point of the new step is not yet a new language, but a sharper compression law: two visible residual presentations should count as the same as soon as they carry the same admissible record. This is the first place where later ZFM-facing or proof-economy layers can aim at strict 1:1 translation without reopening the visible branch picture. In the current 3.5.x reading this also fixes the first status boundary: open work should now be read primarily as quantitative or global-identification work, not as permission to introduce a second visible interface.


\subsection*{3.5 Third extension: record interpreter and exact representability principle}
\addcontentsline{toc}{subsection}{3.5 Third extension: record interpreter and exact representability principle}
Once residual-record morphisms have been fixed, the next step is to say not only when two visible presentations are the same, but how a visible theorem is to be read at record level. The key point is that the stabilized 3.4 package should now admit a faithful interpreter: visible statements are no longer primary objects, but are required to be represented by formulas built from the residual record together with chartwise shell coordinate and parity label. This still does not introduce a separate proof language. It only fixes the exact translation interface that any later compression layer must preserve.

\begin{definition}[Residual-record interpreter]
\label{def:mainline-3502}
Let $B$ be a connected admissible shell base with stabilized record
\[
  \mathfrak R(B)=\bigl(\Phi,\Delta_{\mathrm{diag}},\varepsilon_{\mathrm{par}},\chi_{\mathrm{par}},\widehat\chi_{\mathrm{par}}\bigr).
\]
A \emph{residual-record interpreter} on $B$ is an assignment
\[
  \mathcal I_B : \mathrm{Vis}(B) \longrightarrow \mathrm{Rec}(B)
\]
from visible statements in the stabilized generator-atlas-holonomy sector to record-level formulas such that:
\begin{enumerate}[label=(\roman*), leftmargin=2em]
  \item every local branch law, inverse-shell statement, tunnel-observable relation, shell-consistency statement, parity-obstruction claim, holonomy statement, and empirical loop-readout identity is sent to a formula built only from the record $\mathfrak R(B)$ together with chartwise shell coordinate and parity label;
  \item truth is preserved and reflected: a visible statement holds on $B$ if and only if its interpreted record-level formula holds on the same admissible shell charts and loop data;
  \item the interpretation is natural under residual-record isomorphisms, i.e. for every residual-record isomorphism $F:B_1\to B_2$ one has
  \[
    \mathcal I_{B_1}(F^*\sigma)=F^*\bigl(\mathcal I_{B_2}(\sigma)\bigr)
  \]
  for all visible statements $\sigma$ in the stabilized sector.
\end{enumerate}
\end{definition}

\begin{theorem}[3.5 exact representability principle]
\label{thm:mainline-3502}
Assume the hypotheses of Theorems~\ref{thm:mainline-3500} and~\ref{thm:mainline-3501}. Then the stabilized visible generator-atlas-holonomy sector admits a residual-record interpreter in the sense of Definition~\ref{def:mainline-3502}. More precisely:
\begin{enumerate}[label=(\roman*), leftmargin=2em]
  \item every visible statement in the stabilized sector has a record-level representative built from $\mathfrak R(B)$ together with chartwise shell coordinate and parity label;
  \item this representative is unique up to residual-record isomorphism;
  \item proof transport along residual-record isomorphisms commutes with interpretation;
  \item consequently, the stabilized visible sector is \emph{exactly representable} at record level: later proof compression may replace visible statements by record-level formulas, but may not change their truth content or open a new visible primitive.
\end{enumerate}
\end{theorem}

\begin{proof}
Theorem~\ref{thm:mainline-3500} proves that every visible statement in the stabilized 3.4 handoff sector factors through the record $\mathfrak R(B)$ together with chartwise shell coordinate and parity label. Theorem~\ref{thm:mainline-3501} then shows that this factorization is invariant under residual-record isomorphisms. Therefore one defines $\mathcal I_B(\sigma)$ by replacing each visible generator, shell, diagnostic, and parity-holonomy ingredient of $\sigma$ by its record-level representative. By construction, this translation preserves and reflects truth, because no visible theorem in the stabilized sector depends on any further primitive datum. Naturality under residual-record isomorphisms follows from the faithful translation principle. Uniqueness up to record-isomorphism is immediate, since any two representatives of the same visible statement differ at most by a transport that preserves the entire record package.
\end{proof}

\begin{corollary}[3.5 interpreter corollary: admissible proof interface]
\label{cor:mainline-3502}
Any later proof-economy layer, formal front-end, or set-theoretic presentation intended to compress the stabilized visible branch theory must factor through the residual-record interpreter. In particular, no later compression is admissible unless it is exact at record level.
\end{corollary}

\begin{proof}
This is the direct consequence of Theorem~\ref{thm:mainline-3502}. Once the visible stabilized sector is exactly representable at record level, any faithful later interface must preserve this representation rather than replace it by a different visible primitive basis.
\end{proof}

\paragraph{3.5.2 working remark.}
The present step is still intentionally weaker than a full new language. It only fixes the precise interface that any later language-level simplification must respect. This is already enough to say what a faithful front-end would have to do: not create new visible objects, but interpret the existing theory exactly through the admissible residual record. From this point onward, the live fronts should therefore be read as residual quantitative or global tasks --- for example the continuum-to-measurement correction in $\bsym$, the controlled-emergence problem, the mass hierarchy and generation law, the compact Dixon reformulation, and the remaining Theta-sum balance --- rather than as a search for further visible primitives.



\subsection*{3.6 First extension: conservative frontend principle}
\addcontentsline{toc}{subsection}{3.6 First extension: conservative frontend principle}
The first genuine 3.6 move is not a new visible mechanism, but a theorem saying exactly what a later proof language, ZFM-facing compression layer, or formal frontend is allowed to do. The visible interface fixed in 3.5 may be reorganized, typed, or rendered more economical; it may not be replaced by a second primitive visible basis. This is the point at which the future language project becomes mathematically testable.

\begin{definition}[Conservative frontend over the residual-record interpreter]
\label{def:mainline-3600}
Assume the stabilized generator-atlas-holonomy regime on a connected admissible shell base $B$, and let
\[
  \mathcal I_B : \mathrm{Vis}(B) \longrightarrow \mathrm{Rec}(B)
\]
be the residual-record interpreter of Definition~\ref{def:mainline-3502}. An \emph{admissible conservative frontend} on $B$ consists of
\[
  \mathcal F_B : \mathrm{Vis}(B) \longrightarrow \mathrm{Fr}(B),
  \qquad
  \mathcal J_B : \mathrm{Fr}(B) \longrightarrow \mathrm{Rec}(B),
\]
where $\mathrm{Fr}(B)$ is a formal compression layer, such that:
\begin{enumerate}[label=(\roman*), leftmargin=2em]
  \item \textbf{factorization through the record interpreter:}
  \[
    \mathcal J_B\circ \mathcal F_B = \mathcal I_B;
  \]
  \item \textbf{truth preservation and reflection on the visible image:} for every visible statement $\sigma$, the frontend statement $\mathcal F_B(\sigma)$ holds if and only if the record-level formula $\mathcal J_B(\mathcal F_B(\sigma))$ holds;
  \item \textbf{naturality under residual-record isomorphisms:} for every residual-record isomorphism $G:B_1\to B_2$ there is an induced frontend transport $G^*_{\mathrm{Fr}}$ such that
  \[
    \mathcal F_{B_1}(G^*\sigma)=G^*_{\mathrm{Fr}}\bigl(\mathcal F_{B_2}(\sigma)\bigr)
    \qquad\text{and}\qquad
    \mathcal J_{B_1}\bigl(G^*_{\mathrm{Fr}}(\tau)\bigr)=G^*\bigl(\mathcal J_{B_2}(\tau)\bigr);
  \]
  \item \textbf{no new visible primitives on the visible image:} every primitive symbol needed to express $\mathcal F_B(\sigma)$ for a visible statement $\sigma$ is definable, on the visible image, from the residual-record package $\mathfrak R(B)$ together with chartwise shell coordinate and parity label.
\end{enumerate}
\end{definition}

\begin{theorem}[3.6 conservative frontend principle]
\label{thm:mainline-3600}
Assume the hypotheses of Theorems~\ref{thm:mainline-3500}, \ref{thm:mainline-3501}, and~\ref{thm:mainline-3502}. Let $(\mathrm{Fr}(B),\mathcal F_B,\mathcal J_B)$ be an admissible conservative frontend in the sense of Definition~\ref{def:mainline-3600}. Then the frontend is conservative over the stabilized visible sector. More precisely:
\begin{enumerate}[label=(\roman*), leftmargin=2em]
  \item every visible theorem is represented in the frontend by a statement whose record-level meaning is exactly the one fixed by the residual-record interpreter;
  \item frontend proof transport along residual-record isomorphisms commutes with visible proof transport and with record interpretation;
  \item no frontend proof concerning the stabilized visible sector can depend, on its visible image, on a primitive datum not already definable from $\mathfrak R(B)$ together with chartwise shell coordinate and parity label;
  \item if $(\mathrm{Fr}'(B),\mathcal F'_B,\mathcal J'_B)$ is a second admissible conservative frontend, then the two frontends agree on the visible image up to unique record-level translation. In particular, they are equivalent as compressions of the same stabilized visible theory.
\end{enumerate}
Equivalently, any later admissible language layer is conservative precisely by factoring through the residual-record interpreter and introducing no new visible primitive basis.
\end{theorem}

\begin{proof}
By Theorem~\ref{thm:mainline-3502}, every visible statement in the stabilized sector already has an exact record-level representative, unique up to residual-record isomorphism. Condition~(i) in Definition~\ref{def:mainline-3600} forces every frontend rendering of a visible statement to pass through exactly that representative. Condition~(ii) then implies that frontend truth on the visible image is preserved and reflected by the same record-level formula, so the frontend cannot change the theorem content of the stabilized sector. Condition~(iii) aligns frontend transport with the residual-record transport fixed in Theorem~\ref{thm:mainline-3501}, hence visible proof transport, frontend proof transport, and record interpretation commute. Finally, Condition~(iv) excludes the only remaining failure mode: a formally new symbol that would secretly function as a new visible primitive. Therefore the frontend is conservative over the visible theory. If a second conservative frontend is given, both must factor through the same interpreter and hence agree on the visible image up to the unique record-level identification already fixed by Theorem~\ref{thm:mainline-3502}.
\end{proof}

\begin{corollary}[3.6 admissibility test for later language layers]
\label{cor:mainline-3600}
Any later ZFM-facing layer, typed proof language, computational frontend, or proof-economy syntax is admissible for the stabilized visible sector only if it is introduced as a conservative frontend in the sense of Definition~\ref{def:mainline-3600}. In particular, later language work is no longer judged by expressive novelty, but by exact factorization through the residual-record interpreter.
\end{corollary}

\begin{proof}
This is the direct content of Theorem~\ref{thm:mainline-3600}. Once 3.5 has fixed the visible residual interface and its exact record-level interpretation, the only mathematically faithful later language layers are those that remain conservative over that same interface.
\end{proof}

\begin{lemma}[3.6.1 composition and primitive non-drift]
\label{lem:mainline-3610}
Let
\[
  \mathcal F_B:\mathrm{Vis}(B)\to \mathrm{Fr}(B), \qquad
  \mathcal J_B:\mathrm{Fr}(B)\to \mathrm{Rec}(B)
\]
and
\[
  \mathcal F'_B:\mathrm{Fr}(B)\to \mathrm{Fr}'(B), \qquad
  \mathcal J'_B:\mathrm{Fr}'(B)\to \mathrm{Rec}(B)
\]
be admissible conservative frontends over the same stabilized visible sector, with the second layer conservative over the record semantics already fixed by the first. Then the composite presentation
\[
  \mathcal F'_B\circ \mathcal F_B:\mathrm{Vis}(B)\to \mathrm{Fr}'(B)
\]
is again an admissible conservative frontend. In particular, conservative compression is closed under composition.

Moreover, no primitive drift can occur under admissible conservative extension: every primitive symbol in the downstream frontend \(\mathrm{Fr}'(B)\) that contributes to visible semantics remains definable, on the visible image, from the residual-record package
\[
  \mathfrak R(B)=\bigl(\Phi,\Delta_{\mathrm{diag}},\varepsilon_{\mathrm{par}},\chi_{\mathrm{par}},\widehat\chi_{\mathrm{par}}\bigr)
\]
together with the chartwise shell coordinate and parity label \((s,\pi)\).
\end{lemma}

\begin{proof}
By admissibility of the first layer,
\[
  \mathcal J_B\circ \mathcal F_B=\mathcal I_B.
\]
By admissibility of the second layer over the already fixed record semantics of the first,
\[
  \mathcal J'_B\circ \mathcal F'_B=\mathcal J_B
\]
on the image relevant for visible statements. Therefore
\[
  \mathcal J'_B\circ \mathcal F'_B\circ \mathcal F_B
  =\mathcal J_B\circ \mathcal F_B
  =\mathcal I_B,
\]
so the composite frontend still factors through the residual-record interpreter. Truth preservation and reflection pass through composition, and transport compatibility is inherited because both frontend layers commute with residual-record isomorphisms. Finally, if a downstream primitive were to carry genuinely new visible information, its record image would define a new visible primitive outside \(\mathfrak R(B)\) and \((s,\pi)\), contradicting admissibility. Thus conservative frontends compose and cannot introduce primitive drift.
\end{proof}

\paragraph{3.6.1 working remark.}
The role of this lemma is to prevent later language engineering from re-opening the visible sector by iteration. Once a frontend layer is certified as conservative, any further admissible compression must remain conservative over the same residual-record interface. So 3.6 is now not only an admissibility principle for one future language layer, but a closure principle for chains of later language layers.

\paragraph{Status of the 3.6 line.}
Section~3.6 still does not build the later language itself. Its role is narrower and structural: it fixes the admissibility constraints that any later ZFM-facing layer, typed proof frontend, or compressed proof-economy syntax must satisfy. The visible interface remains the residual-record interface established in Section~3.5; Section~3.6 only certifies which later language layers are conservative over that interface.



\paragraph{3.6.0 working remark.}
This is the first point at which the future language project becomes structurally disciplined enough to be tested. A frontend is now admissible only if it acts as a conservative presentation of the already fixed visible theory. That shifts the burden of later language work completely: the task is no longer to invent a prettier or stronger visible basis, but to prove exact factorization, transport-compatibility, and primitive non-drift over the residual record.

\subsection*{3.7 OP-first consolidation: residual downstream reduction}
\addcontentsline{toc}{subsection}{3.7 OP-first consolidation: residual downstream reduction}
The first 3.7 step is intentionally \emph{not} a new QAM-, ZFM-, or language-layer construction. After the 3.5 residual-record interface and the 3.6 conservative-frontend discipline, the natural next question is whether the remaining historical OPs can now be reduced further \emph{without} enlarging the visible ontology. In the present stabilized reading, the answer is yes: the remaining live fronts are to be treated as residual downstream tasks, not as invitations to reopen the visible basis.

\begin{definition}[Residual downstream OP]
\label{def:mainline-3700}
A historical open problem is called \emph{residual downstream} if its remaining content can be formulated entirely in terms of the residual-record package
\[
  \mathfrak R(B)=\bigl(\Phi,\Delta_{\mathrm{diag}},\varepsilon_{\mathrm{par}},\chi_{\mathrm{par}},\widehat\chi_{\mathrm{par}}\bigr),
\]
the chartwise shell/parity data \((s,\pi)\), and conservative downstream constructions over the same record semantics. Equivalently, solving such a problem may require new estimates, global identifications, asymptotic control, or controlled emergence arguments, but may not require a new visible primitive basis.
\end{definition}

\begin{theorem}[3.7 residual downstream reduction principle]
\label{thm:mainline-3700}
Assume the stabilized visible generator-atlas-holonomy sector, the exact representability result of Section~3.5, and the conservative frontend discipline of Section~3.6. Then the remaining historical OPs split as follows.
\begin{enumerate}[label=(\roman*), leftmargin=2em]
  \item \textbf{Already closed as separate historical OPs.}
  OP~2 as an independent operator problem, OP~3 as a primary open problem, and OP~5 step~(A) no longer remain open fronts in the main line. Their surviving content is either already absorbed into OP~1-type correction work or reduced to the residual part of OP~5.
  \item \textbf{Reduced residual open core of OP~1.}
  The remaining content of OP~1 is not the existence of a visible frame candidate, but the explicit prime-lattice / continuum-limit correction that moves the continuum fixed-point candidate to the measured matter-frame value.
  \item \textbf{Reduced residual open core of OP~5.}
  The remaining content of OP~5 is the quantified Theta-sum / balance problem from step~(B). In particular, OP~5 is no longer a broad visible-sector problem, but a residual quantitative balance problem downstream of the same stabilized record.
  \item \textbf{Other live fronts remain downstream.}
  The field-equation / controlled-emergence problem, the mass-hierarchy problem, the three-generation problem, and the compact Dixon reformulation remain open, but their open content is downstream of the residual-record interface. None of them justifies reopening the visible primitive basis fixed in Sections~3.5--3.6.
  \item \textbf{No remaining OP presently forces a second visible interface.}
  Consequently, the current OP-first program may proceed by reduction, estimation, and controlled downstream closure alone: in the present stabilized reading, any further progress on the remaining live fronts must factor through the existing residual-record semantics and may not appeal to a new visible primitive layer.
\end{enumerate}
\end{theorem}

\begin{proof}
Item~(i) is the status reclassification already recorded in the 3.5.x cleanup line: OP~2 as a separate operator problem is absorbed, OP~3 is no longer carried as a primary main-line open front, and OP~5 step~(A) is already closed while only step~(B) remains. Item~(ii) is exactly the current reading of OP~1 in the synchronized historical-open-problems layer: the broad frame-candidate issue has collapsed to the explicit prime-lattice / continuum-limit correction. Item~(iii) is the corresponding synchronized reading of OP~5, whose remaining content is the Theta-sum / balance step after the earlier renormalization and step~(A) closure. Item~(iv) is the present status of the remaining live fronts, all of which are explicitly recorded as open but downstream structural or quantitative tasks rather than missing visible primitives. Finally, Item~(v) follows from Theorem~\ref{thm:mainline-3502}, Theorem~\ref{thm:mainline-3600}, and Lemma~\ref{lem:mainline-3610}: once visible truth is exactly representable at record level and all admissible later layers are conservative, no unresolved historical OP can legitimately force a second visible interface in the present stabilized reading without violating the already established residual-record discipline.
\end{proof}

\begin{corollary}[3.7 OP-first corollary: defer later language growth]
\label{cor:mainline-3700}
At the present stage, there is no mathematical need to pull the later QAM-, ZFM-, or proof-language layer upward in order to continue the main line. The next admissible 3.7 work may instead focus directly on residual OP closure --- in particular on the OP~1 correction problem and the remaining OP~5 balance step --- while keeping later language engineering deferred.
\end{corollary}

\begin{proof}
This is the direct consequence of Theorem~\ref{thm:mainline-3700}. Since the remaining live fronts are already downstream of the residual-record interface and do not require a new visible primitive basis, the main line may continue by solving those residual tasks first.
\end{proof}

\paragraph{3.7.1 working remark.}
This step deliberately changes the local priority order. The document now records that the next main-line gains should come from solving the residual OP cores whenever possible, not from prematurely lifting the deferred language program. QAM-, ZFM-, or typed-frontend work therefore remains admissible but non-urgent: it can stay deferred until the remaining downstream OPs have either been closed or reduced further.

\paragraph{3.7.1 release note.}
This public release is recorded under Zenodo DOI:
\begin{quote}
\texttt{10.5281/zenodo.19412478}
\end{quote}
It freezes the OP-first 3.7 line without advancing the deferred QAM/language layer.

The present subsection should be read as a release-grade consolidation claim, not as a declaration that every residual OP is already solved. What is claimed is narrower and more structural: the current chain is strong enough that the remaining live fronts can be pursued as downstream residual tasks over the fixed record interface. This is the point that makes a 3.7 release defensible while leaving the actual OP-closure push to the next cycle.

